\documentclass[a4paper]{article}


\usepackage{graphicx}

\usepackage{amsmath,amssymb}

% web referencing
\usepackage[colorlinks=true, linkcolor=blue, urlcolor=blue]{hyperref}
\newcommand{\seehere}[1]{\href{#1}{see here}

\usepackage{minted}
\usepackage{multirow}
\usepackage{float}
\usepackage{todonotes}
\usetikzlibrary{graphs,graphdrawing}
\usegdlibrary{layered}
\usepackage{forest}
\usepackage{xstring}
\input{/home/donkarlo/Dropbox/repo/nd_ai_project/src/nd_ai/knoweldge_representation/language/formal/computer/programming/tex/latex/code_clips/external_dot.tex}
\title{}
\date{}

\begin{document}
\maketitle
"Adaptive Control of Thought—Rational") is a cognitive architecture mainly developed by John Robert Anderson and Christian Lebiere at Carnegie Mellon University. Like any cognitive architecture, ACT-R aims to define the basic and irreducible cognitive and perceptual operations that enable the human mind. In theory, each task that humans can perform should consist of a series of these discrete operations. \seehere{https://en.wikipedia.org/wiki/ACT-R}
% \bibliographystyle{plain}
% \bibliography{/home/donkarlo/Dropbox/repo/shared_research/bibliography/bibliography.bib}
\end{document}